\subsection{Exact amplitude and row-norm comparisons}
\label{sec:flag-exact-norm-comparison}

We keep the actual local data of
Theorem~\ref{thm:flag-local-data}, their original
Harish--Chandra parameter, and every prescribed pinned
quotient measure. The comparison scalar in the row
consumer is allowed to depend on the fixed row. It must
be chosen before all represented points and radial
sections, but need not have a separate uniform bound.
The same scalar occurs in the pointwise factorization
and the radial integral, and can be cancelled exactly.
This observation removes a normalization loss from the
remaining pointwise problem.\cite{current-construction-record}

Here and below \(r\ge0\), \(k\ge1\), \(s=r+1\),
\(n=2r+k+1\), \(m=2r+2\), and \(f\) is the number of
real places. We work in the concrete standing amplitude
regime, with its integrable discrete-series projector,
original arithmetic ties, and large-gap smooth-decay
reserve. We use the period and mixed Jacquet identities
and the all-cell factorization from
Section~\ref{sec:all-cell-pieri-transfer}. The two radial
equations in the given admissible mass datum will be
displayed explicitly below. Thus no construction of a
different mass datum is implicit in the argument.

\begin{theorem}[Exact comparisons for the actual row]
\label{thm:flag-exact-norm-comparison}
Fix the original pair of words \(u,u'\), the actual
matched local data, an admissible mass datum, and an
admitted pinned global quotient measure. Let
\(J_{\rm pin}^2(w)\) denote the squared Hilbert--Schmidt
norm of the literal pinned global orbital operator at
a represented relevant nonzero-Gram row point \(w\).
There are a constant \(K_0>0\) and, at each real place
\(v\), a single finite compact-finite coefficient \(F_v\)
of the original local discrete series such that
\begin{align}
 \mathrm{OrbInt}(F_v,g_v)
       &=\mathcal J_v(g_v),                         \label{eq:flag-exact-local}\\
 jNorm2(F_v)&=M_v<\infty,                            \label{eq:flag-exact-radial}\\
 \prod_v\|\mathcal J_v(g_v)\|_{\rm HS}^{\,2}
       &=K_0J_{\rm pin}^2(w),                       \label{eq:flag-exact-point}\\
 \prod_v M_v&=K_0I_{\rm re}.                        \label{eq:flag-exact-row}
\end{align}
The first equality holds at every relevant nonzero-Gram
cell, for every admitted orbital quotient measure at
canonical root scale \(1\). Here \(\mathcal J_v\) is the
actual normalized local amplitude and \(M_v\) its original
squared radial integral. The second and fourth equalities
hold for every admitted radial section. The constants
\(K_0\) and \(F_v\) are independent of the represented
point, the section, and the auxiliary finite-place data
in the point-sum consumer.

Consequently the actual uniform amplitude and row-norm
comparisons hold with
\[
 \theta=\theta'=1,\qquad K_v=1,\qquad
 C_1=C_a=C_n=0,\qquad C_{\rm fac}=K_0^{-1}.
                                                        \tag{EN1}
\]
In particular
\[
 C_{\rm fac}\prod_vjNorm2(F_v)=I_{\rm re}.
                                                        \tag{EN2}
\]
For the generic pinned consumer take \(u'=u\).
\end{theorem}

\begin{proof}
\emph{A coefficient that matches the actual amplitude
on every cell.}
First work at one real place. Write \(G\) for the
unitary group and \(S=N\rtimes K\), \(S'=N'\rtimes K'\)
for the two original Bessel subgroups, with their fixed
compact and unipotent product Haar measures. The two
period-normalized maps are denoted
\[
 B:\mathcal H^\infty\longrightarrow W,\qquad
 B':\mathcal H^\infty\longrightarrow W'.
\]
The maps are nonzero for these actual data. To see the
point that is needed here, the integrable coefficient
of the given lowest vector defines a convolution
idempotent on the equivariant quotient Hilbert space.
It fixes the nonzero given seed. On compactly supported
smooth sections, disintegration of its kernel and the
original period identity express that kernel as the
rank-one kernel of the period-normalized map. If this
map were zero, the idempotent would vanish on a dense
subspace and hence everywhere, contradicting its action
on the seed. The absolute convergence needed for this
kernel computation is supplied by the original smooth
coefficient decay, as in the extraction of
Section~\ref{sec:flag-local-data} and the quotient-mass
proof in \texttt{global-norm-frontier-quotient-mass.tex}
of the construction record.\cite{current-construction-record}
More explicitly, with inner products linear in the first variable,
the projector coefficient and its quotient kernel are
\[
 \begin{split}
 e(h)&=d_\sigma\langle v_0,\sigma(h)v_0\rangle,\\
 K(x,y)&=\int_S e(x^{-1}ty)\tau(t)\,dt\\
       &=d_\sigma\bigl[B\sigma(x)v_0
                              \otimes(B\sigma(y)v_0)^*\bigr].
 \end{split}
\]
The last equality is the conjugate of the period identity (AP1),
tested on two vectors of \(W\). If \(B=0\), this kernel and hence
the projector vanish, contradicting its nonzero fixed seed.
The same argument applies to \(B'\). This establishes
nonvanishing from the actual seed, rather than from
multiplicity one alone.

Multiplicity one now identifies the actual local
functions with
\[
 \Phi(g)=\alpha B\sigma(g)v_0,\qquad
 \Phi'(g)=\alpha'B'\sigma(g)v_0,
 \qquad \|v_0\|=1,
                                                        \tag{EN3}
\]
where \(\alpha,\alpha'\ne0\) are finite constants.
The vector \(v_0\) is the original common smooth lowest
vector, normalized to length one. Let \(c_N>0\) be the
given scalar in the definition of the actual amplitude:
\[
 \mathcal J(g)=\frac{\alpha\overline{\alpha'}}{c_N}
       \int_{R'\backslash G}
       B\sigma(gx)v_0\otimes(B'\sigma(x)v_0)^*
                    \,d\dot x ,
 \quad R'=g^{-1}(N\cap gN'g^{-1})g .
                                                        \tag{EN4}
\]
Changing the lowest vector to unit length only changes
\(\alpha,\alpha'\). The factor \(c_N\) is retained
literally.

Compact-finite smooth vectors are dense in the smooth
topology, so choose \(x,y\) with \(Bx\ne0\) and
\(B'y\ne0\). Put
\[
 S_0=\|Bx\|^2\|B'y\|^2>0,\qquad
 F_0(h)=\langle\sigma(h)y,x\rangle
                   \bigl(Bx\otimes(B'y)^*\bigr).
\]
The all-cell factorization
Theorem~\ref{thm:all-cell-relative-factorization} gives
\(\mathrm{OrbInt}(F_0,g)=S_0T_g\), with the original
root measure on \(R=N\cap gN'g^{-1}\). Define
\[
 q_{B'}=\int_{S'\backslash G}
                \|B'\sigma(x)v_0\|^2\,d\dot x .
                                                        \tag{EN5}
\]
By (EN3) and the positive finite quotient mass of the
given \(\Phi'\), one has \(0<q_{B'}<\infty\). Keeping this
number unevaluated avoids any convention concerning
compact-type dimension or formal degree.

We claim the exact identity
\[
 \int_{R'\backslash G}
       B\sigma(gx)v_0\otimes(B'\sigma(x)v_0)^*
                         \,d\dot x=q_{B'}T_g .
                                                        \tag{EN6}
\]
Indeed, divide (AP4) by its fixed nonzero first factor.
Its remaining integral, evaluated at a vector \(w\in W'\),
is
\[
 \int_{R'\backslash S'}
 B\sigma(gs')z\,\langle w,\tau'(s')a'\rangle\,d\dot s'
           =\langle B'z,a'\rangle T_gw .
                                                        \tag{EN7}
\]
Set \(z=\sigma(x)v_0\), \(a'=B'\sigma(x)v_0\), and
integrate over \(S'\backslash G\). Equivariance gives
\(\tau'(s')a'=B'\sigma(s'x)v_0\). The two successive
Weil identities disintegrate \(R'\backslash G\) over
\(S'\backslash G\), with fibre \(R'\backslash S'\).
They give (EN6), with the exact original measures.

Here is an absolute-convergence check for that last
operation. It also proves that the amplitude in (EN4)
exists at each cell, before any uniform pointwise
estimate is used. Choose a symmetric submultiplicative
Frobenius norm \(N\ge1\) on \(G\). The original
smooth-decay theorem gives
\[
 |\langle\sigma(h)v_0,z\rangle|
                  \le C_zN(h)^{-a}
                                                        \tag{EN8}
\]
for each fixed smooth test vector \(z\). Choose
\[
 b>64s^2(2n^2+1),\qquad a>b+4n^2.
                                                        \tag{EN9}
\]
For instance \(b=256s^2(n^2+1)\) is sufficient, and
the original large-gap reserve provides a larger \(a\).
The group-volume bound in
the quotient-mass proof in the construction record gives
\(\int_G N(x)^{b-a}\,dx<\infty\).
For any \(h_1,h_2\in G\), symmetry and
submultiplicativity imply
\[
 \begin{split}
 &\int_GN(h_1x)^{-a}N(h_2x)^{-a}\,dx\\
 &\hspace{5mm}\le
 N(h_1h_2^{-1})^{-b}
                 \int_GN(y)^{b-a}\,dy .
 \end{split}                                           \tag{EN10}
\]
Indeed \(N(h_1h_2^{-1})\le N(h_1x)N(h_2x)\);
allocate \(b\) powers to this inequality, discard
one remaining factor, which is at most one, and
apply left invariance of Haar measure.

Expand both period maps through the period identity,
using fixed test vectors on which they are nonzero.
For the resulting three-variable integrand the
simultaneous \(R'\)-action is
\[
 a'\cdot(s,s',x)
       =((ga'g^{-1})s,\ a's',\ a'x).
\]
Relevance cancels the two character factors.
The quotient of \(S\times S'\times G\) by this action
has two disintegrations: one first integrates
the two full period subgroups over \(R'\backslash G\),
and the other first integrates \(G\) over
\(\Delta_gR\backslash(S\times S')\). Both follow by
composing the defining Weil identities. Tonelli's
theorem applies to the absolute values. Equation (EN10),
with \(h_1=s^{-1}g\), \(h_2=s'^{-1}\), bounds the entire
absolute integral by a fixed finite constant times
\[
 \left(\int_GN(y)^{b-a}\,dy\right)
 \left(\int_{\Delta_gR\backslash(S\times S')}
                    N(s^{-1}gs')^{-b}\,d\dot{(s,s')}\right).
                                                        \tag{EN11}
\]
The second integral is finite by the terminal quotient
properness estimate used in the all-cell factorization:
its real noncompact dimension is at most \(2n^2\) and
its inverse-coordinate exponent is at most \(64s^2\).
Thus (EN9) suffices. This proves the required absolute
integrability and justifies Fubini in (EN6). The
properness constant may depend on this fixed \(g\);
it enters no uniform inequality.

It now suffices to take the single coefficient
\[
 F=\frac{\alpha\overline{\alpha'}q_{B'}}{c_NS_0}F_0 .
                                                        \tag{EN12}
\]
Equations (EN4), (EN6) and the all-cell factorization
prove \eqref{eq:flag-exact-local} on every cell. The
coefficient and its scalar have been chosen before
\(g\) and before the orbital quotient measure. At
canonical root scale, changing the measure representing
the quotient does not change these integrals: the Weil
identity determines its pushforward to the quotient,
and the integrands are invariant along the stabilizer.

On the big cell the equality just proved identifies
the radial integral defining \(jNorm2(F)\) with the
original \(M_v\), for the same section and the same
weight. The mixed Jacquet identity (AP2) gives
\[
 jNorm2(F)=c_{\rm rad}d_\sigma^{-1}|s_\sigma(F)|^2
                                                           <\infty .
                                                        \tag{EN13}
\]
This proves \eqref{eq:flag-exact-radial} without a
pointwise bound. No estimate for the scalar in (EN12)
is required.

\emph{The scalar in the actual pinned factorization.}
Write \(V_K>0\) for the original level volume and put
\[
 c_u=\frac{\mu_{N,u}(F_{N,u})\mu_{H,u}(H_u)}
                              {\dim V_u},
 \qquad c_{0X,u}=V_Kc_u .
                                                        \tag{EN14}
\]
The mass datum uses \(c_u\); the Poincare normalization
uses \(c_{0X,u}\). For each place let \(L_{u,v}\) be the
original positive finite squared quotient mass and put
\[
 c_{u,v}=\frac{L_{u,v}}{\dim V_{u,v}}>0,\qquad
 \kappa_v=(c_{u,v}c_{u',v})^{-1}.
                                                        \tag{EN15}
\]
All displayed dimensions are positive, since the
corresponding seed is nonzero. The literal definitions
of the matched local data give
\[
 \mathcal J_v
       =a_vU_vJ_v(U'_v)^{-1},\qquad
 a_v=c_{u,v}^{-1}\sqrt{\frac{c_{u,v}}{c_{u',v}}},
 \qquad |a_v|^2=\kappa_v ,
                                                        \tag{EN16}
\]
where \(U_v,U'_v\) are the unitary frame and compact-type
identifications and \(J_v\) is the original unnormalized
local operator. The squared global scalar is exactly
\[
 s_{\rm gl}
       =\left(\frac{V_K}{\sqrt{c_{0X,u}c_{0X,u'}}}\right)^2
       =(c_uc_{u'})^{-1}.
                                                        \tag{EN17}
\]
Thus set
\[
             K_0=c_uc_{u'}\prod_v\kappa_v>0 .
                                                        \tag{EN18}
\]

The canonical root measures in (EN16) are exactly the
transported original ones. In root coordinates the
conjugating frame map is a permutation and reindexing
with scale one. In the original character convention,
the density per complex root is
\(\kappa_{\rm char}/\pi\), where
\(\kappa_{\rm char}=\mathrm{ka}/g_{\rm ar}\);
the transported formula
\((2\,\mathrm{ka}/g_{\rm ar})/(2\pi)\) is the same
density. Here \(g_{\rm ar}>0\) is the original
arithmetic character scale, not the orbital
representative. There is no extra root-scale factor.
The triangular root-coordinate group law preserves
this Haar measure, including in cells with vanishing
internal minors. Relevance makes the rank-one
integrand invariant along the stabilizer, so the
preceding Weil-uniqueness argument applies to every
admitted pinned quotient measure.

The product Weil formula for the original global
quotient and the tensor-product identity for the
Hilbert--Schmidt norm now give
\[
 J_{\rm pin}^2(w)=s_{\rm gl}\prod_v\|J_v(g_v)\|_{\rm HS}^2,
 \qquad
 \prod_v\|\mathcal J_v(g_v)\|_{\rm HS}^2
             =\left(\prod_v\kappa_v\right)
                         \prod_v\|J_v(g_v)\|_{\rm HS}^2 .
\]
These are \eqref{eq:flag-exact-point} and (EN18).
No division by a pointwise orbital value has been
made, so zero values are included.

\emph{The same scalar in the radial integral.}
Put \(\rho=\mathrm{rd}(E)>0\). For \(0\le a<r\), set
\[
       \lambda_a=\rho^{-2(r-a)},\qquad
       M_a=2(r-a)+k+1 .
\]
The original diagonal Gram parameter \(z\) and the
normalized local radial parameter \(q\) satisfy
\[
       q_a=\lambda_a z_a\quad(a<r),\qquad q_r=z_r .
                                                        \tag{EN19}
\]
The original determinant and diagonal Jacobian are
\[
 \begin{split}
 G(\operatorname{diag}z)
    &=\prod_{a<r}|z_a|^2\,(|z_r|^2-TT'),\\
 \operatorname{Jac}(\operatorname{diag}z)
    &=\prod_{a<r}|z_a|^{4(r-a)}.
 \end{split}
\]
Consequently their product weight is
\[
 W(z)=\prod_{a<r}|z_a|^{2(M_a-2)}
                         (|z_r|^2-TT')^{k-1}.
                                                        \tag{EN20}
\]
The arithmetic parameter measure is \(2\,dx\,dy\)
in each complex coordinate, whereas \(M_v\) uses
\(dx\,dy\). The real Jacobian of (EN19) is
\(\prod_{a<r}\lambda_a^2\). It follows that
\[
 K_{\rm pl}=2^{r+1}
              \prod_{a<r}\lambda_a^{-2(M_a-1)}>0,
 \qquad
       \left(\prod_v\kappa_v\right)I_C
                        =K_{\rm pl}^{\,f}\prod_vM_v .
                                                        \tag{EN21}
\]
First interpret this equation as an equality of
nonnegative extended integrals. It follows by Tonelli
and the displayed change of variables, with the exact
Hilbert--Schmidt scalar in (EN16). Two representatives
with the same diagonal Gram data differ by the given
two Bessel subgroups; their amplitudes differ by
unitary factors. Hence the calculation holds for
every admitted section. The excluded hyperplanes
\(z_a=0\) and the boundary circle
\(|z_r|^2=TT'\) have Lebesgue measure zero. This remains
valid for \(k=1\), when the boundary weight has exponent
zero, and for \(r=0\), when the empty product gives
\(K_{\rm pl}=2\).

By \eqref{eq:flag-exact-radial} the right side of
(EN21) is finite. The positive finite factor on its
left proves finiteness of the original carrier energy.
Only now pass to real values; in particular finiteness
has not been inferred from a real-valued truncation
of an extended integral.

The elementary exponent identity
\[
 \sum_{j=1}^r4j(2j+k)
   =2r(r+1)(k-1)+\frac{r(r+1)(8r+10)}3
\]
gives
\[
 K_{\rm pl}^{\,f}
   =2^{f(r+1)}
       \rho^{2r(r+1)(k-1)f}
       \rho^{fr(r+1)(8r+10)/3}.                         \tag{EN22}
\]
The two literal radial equations in the admissible
mass datum are
\[
 \begin{split}
 (c_uc_{u'})^{-1}I_C&=2^{f(r+1)}I_{\rm rat},\\
 I_{\rm rat}
    &=\rho^{2r(r+1)(k-1)f}
           \rho^{fr(r+1)(8r+10)/3}I_{\rm re}.
 \end{split}                                           \tag{EN23}
\]
Combining (EN21)--(EN23) and cancelling the positive
finite \(K_{\rm pl}^{\,f}\) proves
\[
      \prod_vM_v
        =c_uc_{u'}\left(\prod_v\kappa_v\right)I_{\rm re}
        =K_0I_{\rm re},
\]
which is \eqref{eq:flag-exact-row}.

\emph{The consumer quantifiers.}
The matched data have canonical root scale \(1\);
at \(C_1=0\) the permitted root window consists of
that scale alone. Equation \eqref{eq:flag-exact-local}
therefore supplies the amplitude comparison with
\(K_v=1\) and \(C_a=0\), for all admitted quotient
measures. Equations \eqref{eq:flag-exact-point} and
\eqref{eq:flag-exact-row} supply the pointwise product
factorization and row-norm comparison with
\(C_{\rm fac}=K_0^{-1}\), \(C_n=0\).
The row consumer chooses \(C_{\rm fac}\) after fixing
the row but before all represented points, all
sections, and its auxiliary \(\nu_L,n_{\rm gen}\);
(EN18) has exactly this dependence. The local
coefficient scalars, local masses, and \(K_0\) can
vary with the arithmetic data, but cancel before
any uniform inequality is applied. This proves
(EN1)--(EN2).
\end{proof}

\paragraph{Boundary of the comparison theorem.}
The theorem supplies the uniform amplitude and
row-norm part of the actual generic pinned route.
Combined with the pointwise estimate (AP6) and the
weighted rotated summation of
Theorem~\ref{thm:rotated-sum-from-near-bp}, it gives the
original \(\mathrm{StepThmBPinnedG}\) point sum with its
original \(I_{\rm re}\) and pinned measure. There is no
additional exponent from amplitude normalization or
row norm. The pointwise estimate for the required family is
proved in Theorem~\ref{thm:archimedean-pointwise-sufficient},
with its full reserve checked in
Proposition~\ref{prop:archimedean-ap7-reserve}.
This cancellation does not identify the retained
theta columns or establish the positive-trace
concentration required later.
